\documentclass[11pt]{article}
\usepackage[margin=0.75in]{geometry}
\usepackage{lmodern}
\usepackage{microtype}
\usepackage{enumitem}
\usepackage{listings}
\usepackage[T1]{fontenc}

\lstset{
  basicstyle=\ttfamily\small,
  columns=fullflexible,
  keepspaces=true,
  showstringspaces=false,
  breaklines=true,
  frame=single,
  xleftmargin=0.4em,
  xrightmargin=0.4em,
  aboveskip=0.55em,
  belowskip=0.55em
}

\pagestyle{plain}
\setlength{\parskip}{0.35em}
\setlength{\parindent}{0pt}
\setlist[enumerate]{leftmargin=1.45em, itemsep=0.7em, topsep=0.35em}
\setlist[itemize]{leftmargin=1.25em, itemsep=0.2em, topsep=0.25em}

\begin{document}

\begin{center}
{\LARGE\bfseries Midterm 1}\\[0.3em]
DATA 412/612 \quad Statistical Programming in R\\[0.15em]
October 14, 2026 \quad 75 minutes \quad 100 points
\end{center}

Open \texttt{midterm-1.qmd}. Put your name in the subtitle.
That file already has a heading and a code chunk for each part. Write your answers there.
Code is already set to show in the rendered file.
Render to HTML (the default). PDF is fine too. Submit the \texttt{.qmd} and the rendered file before the 75 minutes are up.

\begin{itemize}
  \item You may use your notes, the R help files, and the internet.
  \item You may not communicate with anyone else.
  \item Use the tidyverse where a tidyverse verb exists. Question 1 is a function and a pipe; base R is appropriate there.
  \item \texttt{library(tidyverse)} is enough. \texttt{mpg} and \texttt{diamonds} are in ggplot2. \texttt{starwars} is in dplyr. \texttt{mtcars} is in base R.
  \item Either pipe is fine: \texttt{|>} or \texttt{\%>\%}.
  \item Partial credit is given for code that runs. If you are stuck, go to the next part.
\end{itemize}

The suggested times are a guide. You have 75 minutes.

\begin{enumerate}

\item \textbf{A function and a pipe} (20 points, about 14 minutes)

\begin{enumerate}[label=\alph*., leftmargin=1.7em, itemsep=0.55em]
  \item (14 points) Write a function \texttt{finite\_ratio(x, y)}.
  \begin{itemize}
    \item If \texttt{x} or \texttt{y} is not numeric, or the two lengths differ, stop with the message \texttt{x and y must be numeric and the same length}.
    \item Drop every position where \texttt{x} or \texttt{y} is not finite. \texttt{NA}, \texttt{NaN}, \texttt{Inf}, and \texttt{-Inf} are not finite.
    \item If fewer than 2 positions remain, return \texttt{NA\_real\_}.
    \item If the mean of the remaining \texttt{y} is 0, return \texttt{Inf}.
    \item Otherwise return the mean of the remaining \texttt{x} divided by the mean of the remaining \texttt{y}. Use a pipe for that ratio.
  \end{itemize}
  Run it on these four inputs. The chunk in \texttt{midterm-1.qmd} already has \texttt{\#| error: true}, so the last call can stop and the document still renders. Leave that error visible.
\begin{lstlisting}
finite_ratio(c(1, NA, 4, Inf), c(2, 3, 6, 1))
finite_ratio(c(1, 2), c(0, 0))
finite_ratio(c(1, NA_real_), c(NA_real_, 4))
finite_ratio(c(1, 2), "a")
\end{lstlisting}

  \item (6 points) Using one pipe and no intermediate objects, start from \texttt{c(2, 4, 8, NA, 16)}, drop the missing value, take the natural log, take the mean, then exp. Report the number. In one sentence, say what goes wrong if the \texttt{NA} is still in the vector when you call \texttt{log()}.
\end{enumerate}

\item \textbf{Rows, columns, and groups} (20 points, about 15 minutes)

Use \texttt{mpg}. Answer each part with one pipeline.

\begin{enumerate}[label=\alph*., leftmargin=1.7em, itemsep=0.55em]
  \item (8 points) Keep cars from \texttt{year} 2008 in class \texttt{compact} or \texttt{midsize} with \texttt{hwy} at least 30. How many cars is that? Show the five with the highest \texttt{hwy}. Break remaining ties by \texttt{cty}, highest first. Keep only \texttt{manufacturer}, \texttt{model}, \texttt{year}, \texttt{cyl}, \texttt{cty}, and \texttt{hwy}.
  \item (7 points) For each \texttt{class}, report the number of cars, the mean of \texttt{hwy}, and the median of \texttt{cty}. Arrange by mean highway mileage, highest first.
  \item (5 points) Add a column \texttt{economy} equal to \texttt{(cty + hwy) / 2}. Place it immediately after \texttt{model}. Drop \texttt{trans} and \texttt{fl}. Show the column names of the result, in order.
\end{enumerate}

\item \textbf{\texttt{case\_when}, \texttt{across}, and distinct rows} (20 points, about 16 minutes)

\begin{enumerate}[label=\alph*., leftmargin=1.7em, itemsep=0.55em]
  \item (8 points) Add a \texttt{size} column to \texttt{starwars} with \texttt{case\_when()}, in this order:
  \begin{itemize}
    \item \texttt{"missing"} if \texttt{height} is missing
    \item \texttt{"large"} if \texttt{height} is at least 200 or \texttt{mass} is at least 200
    \item \texttt{"droid"} if \texttt{species} is \texttt{"Droid"}
    \item \texttt{"other"} otherwise
  \end{itemize}
  Report how many characters fall in each size, largest count first.
  In one sentence, say what would change about the droid count if \texttt{"large"} were not checked before \texttt{"droid"}. Name the character.

  \item (7 points) Drop characters whose \texttt{species} is missing. Keep species that appear more than once. For each of those species, use \texttt{across()} for the median \texttt{height} and the median \texttt{mass} (ignore missing values inside the medians), and report \texttt{n()} as well. \texttt{n()} counts every row in the species, including rows whose height or mass is missing. Arrange by \texttt{n()}, largest first.

  \item (5 points) Two short pipelines.
  \begin{itemize}
    \item Drop \texttt{starwars} rows where \texttt{hair\_color} or \texttt{skin\_color} is missing, then report how many distinct pairs of those two columns remain.
    \item \texttt{mtcars} stores the car name as a row name. Make that a column called \texttt{car}. Keep cars with \texttt{mpg} of at least 30. Return \texttt{car}, \texttt{mpg}, and \texttt{cyl}, arranged by \texttt{mpg}, highest first.
  \end{itemize}
\end{enumerate}

\item \textbf{Pictures and what they say} (20 points, about 15 minutes)

\begin{enumerate}[label=\alph*., leftmargin=1.7em, itemsep=0.55em]
  \item (6 points) From \texttt{diamonds}, plot a histogram of \texttt{carat} with 40 bins and \texttt{theme\_bw()}. Label the horizontal axis \texttt{Carat}. In one sentence, describe the shape. ``It is a histogram of carat'' is not a description.
  \item (10 points) From \texttt{mpg}, plot highway mileage against engine displacement.
  \begin{itemize}
    \item Points: map \texttt{class} to color inside \texttt{geom\_point()} only, and set an alpha so overlapping points stay visible.
    \item One straight smooth in black, dashed, with no gray band. The smooth is not colored by \texttt{class}.
    \item Facet by \texttt{drv}.
    \item Use \texttt{theme\_bw()}. Label the axes \texttt{Displacement (litres)} and \texttt{Highway mpg}.
  \end{itemize}
  \item (4 points) In two sentences, say what happens to highway mileage as displacement increases, and whether that pattern shows up in every \texttt{drv} panel. Name one difference you can see between panels.
\end{enumerate}

\item \textbf{Parse a table, then \texttt{across()}} (20 points, about 10 minutes)

\texttt{halls} is already in \texttt{midterm-1.qmd}. Do not read a file.

\begin{lstlisting}
halls <- tibble(
  name = c("Maple Hall", "Oak Court", "Pine Lofts", "Elm Annex", "Cedar Row"),
  opened = c("March 4 2018", "July 1 2019", "August 15 2016", NA, "January 9 2020"),
  rent = c("$1,250", "980", "$2,100.50", "missing", "$875"),
  units = c("40", "25", "12", NA, "18"),
  occupied = c("TRUE", "FALSE", "TRUE", "N-A", "TRUE")
)
\end{lstlisting}

\begin{enumerate}[label=\alph*., leftmargin=1.7em, itemsep=0.55em]
  \item (12 points) Clean the table. \texttt{opened} is a date written like \texttt{March 4 2018}. \texttt{rent} is a number even when it contains \texttt{\$} or a comma. \texttt{units} is an integer. \texttt{occupied} is logical. Treat \texttt{"missing"} and \texttt{"N-A"} as missing. Print the cleaned table.
  \item (8 points) With one \texttt{across()}, report how many missing values are in each column of the cleaned table.
\end{enumerate}

\end{enumerate}

\end{document}
